\sectionquestion{What makes code complex?}
\section[Complexity]{Complexity}

\begin{frame}{What makes code hard to work with?}
    \source{ousterhout2018philosophy}

    \begin{center}
    \large\textbf{\textcolor{talkaccent}{Complexity}}: anything that makes code hard to
    understand or change
    \end{center}

    \vfill

    \begin{columns}[T]
        \column{0.32\textwidth}
        \centering
        \onslide<2->
        % Working memory holds about four names; the rest spill out
        \begin{tikzpicture}[
            var/.style={font=\tiny\ttfamily, fill=codebg, draw=talkblue, text=talkfg, inner sep=1.5pt},
            spill/.style={var, draw=talkbad, text=talkbad}]
            \useasboundingbox (-0.1,-0.45) rectangle (3.4,1.6);
            \draw[talkdim, rounded corners] (0,0) rectangle (2.4,1.1);
            \node[font=\tiny, text=talkdim] at (1.2,-0.25) {your working memory};
            \node[var] at (0.6,0.8) {mode};
            \node[var] at (1.7,0.8) {scale};
            \node[var] at (0.6,0.3) {offset};
            \node[var] at (1.7,0.3) {done};
            \node[spill] at (2.0,1.4) {threshold};
            \node[spill] at (2.8,0.85) {k};
            \node[spill] at (2.9,0.3) {use\_bg};
        \end{tikzpicture}\par
        \vspace{0.6em}
        {\large\bfseries\textcolor{talkaccent}{Cognitive load}}\par
        \vspace{0.3em}
        {\small\textcolor{talkdim}{Too much to keep in mind at once}}

        \column{0.32\textwidth}
        \centering
        \onslide<3->
        % One value changes: every file that copied it needs an edit
        \begin{tikzpicture}[
            file/.style={font=\tiny\ttfamily, fill=codebg, draw=talkblue, text=talkfg, inner sep=2pt}]
            \useasboundingbox (-0.1,-0.45) rectangle (3.4,1.6);
            \node[font=\scriptsize\ttfamily, draw=talkbad, text=talkbad, inner sep=2pt] (c)
                at (0.45,0.55) {g = 9.81};
            \node[file] (f1) at (2.6,1.25) {forces.py};
            \node[file] (f2) at (2.6,0.55) {energy.py};
            \node[file] (f3) at (2.6,-0.15) {plots.py};
            \foreach \f in {f1, f2, f3}
                \draw[->, talkbad] (c.east) -- (\f.west);
        \end{tikzpicture}\par
        \vspace{0.6em}
        {\large\bfseries\textcolor{talkaccent}{Change amplification}}\par
        \vspace{0.3em}
        {\small\textcolor{talkdim}{One change, edits everywhere}}

        \column{0.32\textwidth}
        \centering
        \onslide<4->
        % Parts you can see, and one you don't know exists
        \begin{tikzpicture}
            \useasboundingbox (-0.1,-0.45) rectangle (3.4,1.6);
            \foreach \i in {0,1,2} \foreach \j in {0,1}
                \draw[talkblue, thick] ({0.2+\i*0.65},{0.05+\j*0.6}) rectangle ++(0.45,0.4);
            \draw[talkbad, thick, dashed] (2.15,0.35) rectangle ++(0.45,0.4);
            \node[text=talkbad, font=\bfseries] at (2.375,0.55) {?};
        \end{tikzpicture}\par
        \vspace{0.6em}
        {\large\bfseries\textcolor{talkaccent}{Unknown unknowns}}\par
        \vspace{0.3em}
        {\small\textcolor{talkdim}{You don't know what you need to know}}
    \end{columns}
    \onslide

    \vfill

\end{frame}




\begin{frame}{Where does complexity come from?}

    \vfill

    \begin{columns}[T]
        \column{0.48\textwidth}
        \centering
        % One function quietly relies on four things outside it
        \begin{tikzpicture}[
            dep/.style={font=\tiny, fill=codebg, draw=talkbad, text=talkfg,
                        inner sep=2pt, align=center}]
            \useasboundingbox (-2.6,-1.3) rectangle (2.6,1.3);
            \node[font=\scriptsize\ttfamily, fill=codebg, draw=talkblue, thick,
                  text=talkfg, inner sep=3pt] (f) at (0,0) {analyse()};
            \node[dep] (a) at (-1.7,0.9) {global variables};
            \node[dep] (b) at (1.7,0.9) {file format};
            \node[dep] (c) at (-1.7,-0.9) {call \texttt{load()} first};
            \node[dep] (d) at (1.7,-0.9) {another object's data};
            \foreach \n in {a, b, c, d}
                \draw[->, talkbad] (f) -- (\n);
        \end{tikzpicture}\par
        \vspace{0.6em}
        {\large\bfseries\textcolor{talkaccent}{Dependencies}}\par
        \vspace{0.3em}
        {\small\textcolor{talkdim}{Code you can't understand on its own}}

        \column{0.48\textwidth}
        \centering
        \onslide<2->
        % One line, three names that hide their meaning
        \begin{tikzpicture}[
            tok/.style={font=\scriptsize\ttfamily, text=talkfg, inner sep=1pt},
            hid/.style={tok, text=talkbad},
            q/.style={font=\small\bfseries, text=talkbad},
            why/.style={font=\tiny, text=talkdim}]
            \useasboundingbox (-2.6,-1.3) rectangle (2.6,1.3);
            \node[tok, anchor=west] (t1) at (-1.9,0) {r =};
            \node[hid, right=6pt of t1] (t2) {data};
            \node[tok, right=6pt of t2] (t3) {*};
            \node[hid, right=6pt of t3] (t4) {0.0259};
            \node[tok, right=6pt of t4] (t5) {+};
            \node[hid, right=6pt of t5] (t6) {tmp};
            \foreach \t in {t2, t4, t6}
                \node[q, above=3pt of \t] {?};
            \node[why, below=4pt of t2] {vague name};
            \node[why, below=16pt of t4] {magic number};
            \node[why, below=4pt of t6] {unclear name};
        \end{tikzpicture}\par
        \vspace{0.6em}
        {\large\bfseries\textcolor{talkaccent}{Obscurity}}\par
        \vspace{0.3em}
        {\small\textcolor{talkdim}{Information that isn't obvious}}
    \end{columns}
    \onslide

    \vfill
    \onslide<3->

    \takeaway{Everything that follows reduces dependencies or obscurity.}

\end{frame}




